\usepackage[spanish]{babel} 
\usepackage[utf8]{inputenc}
%\usepackage[T1]{fontenc} %Codificacion de fuente
\usepackage{comment}
\usepackage{graphicx}
\usepackage{amsmath}
\usepackage{amssymb}
\usepackage{listings}
\usepackage{float}
\usepackage{hyperref}
\usepackage{enumitem}
\usepackage{multirow}
\usepackage{multicol}
\usepackage[table,xcdraw]{xcolor}
\usepackage{url}
\usepackage{karnaugh-map}
\usepackage{longtable}
\usepackage{csquotes}
\usepackage{marvosym}
\usepackage[normalem]{ulem}
\usepackage{array}
\usepackage{booktabs}
\usepackage{cleveref}
\usepackage{algorithmic}
%\usepackage{graphics}
\usepackage{lineno}
\usepackage{ragged2e}
\usepackage{setspace}
\usepackage{lipsum}
\usepackage{threeparttable}
\usepackage{tikz}


\usetikzlibrary{matrix,calc}


\hyphenation{in-te-rrup-to-res}


\hypersetup{
    colorlinks=false,
    hyperindex=true,
    breaklinks=true,
    %linkcolor=blue,
    %filecolor=magenta,      
    %urlcolor=cyan,
    %citecolor=blue,
    pdftitle={Trabajo UNAD},
    pdfpagemode=FullScreen,
}


\selectlanguage{spanish}

% Colores estilo "Dark Mode" suave o editor moderno
\definecolor{codegreen}{rgb}{0,0.6,0}
\definecolor{codegray}{rgb}{0.5,0.5,0.5}
\definecolor{codepurple}{rgb}{0.58,0,0.82}
\definecolor{arduino_teal}{rgb}{0.0, 0.59, 0.62}
\definecolor{backcolour}{rgb}{0.95,0.95,0.92}

% ESTILO BASE (Con soporte de tildes incluido)
\lstdefinestyle{base_style}{
    backgroundcolor=\color{backcolour},   
    numberstyle=\tiny\color{codegray},
    breaklines=true,                 
    numbers=left,                    
    numbersep=7pt,                  
    showstringspaces=false,
    frame=single,
    rulecolor=\color{black!10},
    literate={á}{{\'a}}1 {é}{{\'e}}1 {í}{{\'i}}1 {ó}{{\'o}}1 {ú}{{\'u}}1 {Á}{{\'A}}1 {É}{{\'E}}1 {Í}{{\'I}}1 {Ó}{{\'O}}1 {Ú}{{\'U}}1 {ñ}{{\~n}}1 {Ñ}{{\~N}}1 % Soporte español
}

% ESTILOS ESPECÍFICOS
\lstdefinestyle{python_style}{
    style=base_style,
    language=Python,
    commentstyle=\color{codegreen},
    keywordstyle=\color{magenta},
    stringstyle=\color{codepurple}
}

\lstdefinestyle{arduino_style}{
    style=base_style,
    language=C++,
    commentstyle=\color{gray},
    keywordstyle=\color{arduino_teal}\bfseries,
    stringstyle=\color{blue}
}

\lstdefinestyle{bash_style}{
    style=base_style,
    language=bash,
    basicstyle=\ttfamily\footnotesize\color{white},
    backgroundcolor=\color{black!85},
    keywordstyle=\color{cyan},
    numbers=none
}

\makeatletter
\ifx\estilo\apa
	\usepackage{caption}
    % 1. Definimos un formato de caption con interlineado APA (1.5 o 2.0)
    \DeclareCaptionFormat{apa-listings}{
        \textbf{#1#2} \\%[-0.3cm]
        \textit{#3}
    }
    \captionsetup[lstlisting]{
        format=apa-listings,
        font={stretch=1.0},
        skip=13pt,
        labelsep=none,
        singlelinecheck=false,
        justification=raggedright
    }
    \lstset{
        style=base_style,
        captionpos=t, 
        basicstyle=\ttfamily\small,
        lineskip=0pt,
        linewidth=1.0\linewidth,
        xleftmargin=0.04\linewidth,
    }
\else
	% Configuración para IEEE (Espacio sencillo y abajo)
    %\captionsetup[lstlisting]{
    %    font={small, stretch=1.0},
    %    skip=5pt,
    %    %justification=centering,
    %    labelsep=period,
    %}
    \lstset{
        style=base_style,
        captionpos=b, 
        basicstyle=\ttfamily\footnotesize,
        lineskip=0pt,
        linewidth=1.0\linewidth,
        xleftmargin=0.05\linewidth,
    }
\fi
\makeatother


\providecommand{\inputpath}{./src/} % Carpeta sugerida para tus códigos
\DeclareGraphicsExtensions{.pdf,.jpeg,.png,.eps}

\graphicspath{ {./graphics/} {./img/} {./diagramas/} }

%\makeatletter
%\ifx\estilo\apa
    % Configuración para APA7
%    \renewcommand{\tablename}{Tabla}
%    \renewcommand{\figurename}{Figura}
%	\renewcommand\contentsname{\largeÍndice}
%	\renewcommand{\listfigurename}{\largeÍndice de Figuras}
%	\renewcommand{\listtablename}{\largeÍndice de Tablas}
%	\renewcommand{\lstlistlistingname}{\largeÍndice de Códigos}
%\else
    % Configuración para IEEEtran
    % IEEE usa \footnotesize y Versalitas por defecto, 
    % esto asegura que se mantenga el estilo pero cambie el nombre.
    %\renewcommand{\tablename}{Tabla}
    %\renewcommand{\figurename}{Figura}
    % Si el error persiste en IEEE, forzamos la redefinición de la macro de Babel
    %\addto\captionsspanish{%
    %  \def\tablename{Tabla}%
    %  \def\figurename{Figura}%
    %}
%\fi
%\makeatother

\renewcommand{\tablename}{Tabla}
\renewcommand{\figurename}{Figura}
\renewcommand{\refname}{Referencias Bibliográficas}
\renewcommand{\lstlistingname}{Código}

\ifx\estilo\apa
    \par\medskip
\else
    \renewcommand{\IEEEkeywordsname}{Palabras Clave}
\fi


% --- ENTORNO PARA OBJETIVOS DINÁMICOS ---
\newenvironment{misObjetivosEspecificos}{
    \ifx\estilo\apa
        % Si es APA, no hace nada especial (párrafos normales)
        \par\medskip
    \else
        % Si es IEEE, inicia una lista de puntos
        \begin{itemize}[leftmargin=*]
    \fi
}{
    \ifx\estilo\apa
        \par\medskip
    \else
        \end{itemize}
    \fi
}

% Comando para los ítems que se adapta
\newcommand{\objItem}{
    \ifx\estilo\apa 
        \par \space 
    \else 
        \item 
    \fi
}
